\section{Independent Poisson letter counts}
\label{sec:poisson_independent_counts}

Let $I$ be a finite alphabet, $N=|I|$, and $t\ge0$.  The word law $\mu_t$
assigns mass $e^{-Nt}t^m/m!$ to each word of length $m$.  For
$n\in\N^I$, let $C_n$ be the set of words with letter counts $n$, and put
$m=\sum_{i\in I}n_i$.  The following calculation identifies the fixed-time
count law used in \emph{A two-dimensional area law from a global spectral gap},
\cite{OpenAI2026AreaLaw}, \texttt{09-amplification.tex}, lines 49--54 and
237--253.

\begin{lemma}[Mass of a prescribed count vector]
    \label{lem:poisson_count_event_mass}
    \lean{PoissonWord.measure_countEvent,
        PoissonWord.card_countFiber_mul_weight,
        PoissonWord.measure_countEvent_eq_prod}
    \leanok
    \uses{def:poisson_word_measure, def:poisson_word_count_vector}
    For every $n\in\N^I$,
    \begin{align}
        \mu_t(C_n)
        &= |C_n|\frac{e^{-Nt}t^m}{m!}
         = \prod_{i\in I}\frac{e^{-t}t^{n_i}}{n_i!}.
        \label{eq:poisson_count_event_mass}
    \end{align}
\end{lemma}

\begin{proof}
    \leanok
    \uses{lem:poisson_word_count_fiber}
    Every word in $C_n$ has length $m$.  Sum its common singleton mass
    over the finite set $C_n$.  The multinomial identity
    $|C_n|\prod_i n_i!=m!$ cancels the factorial in the word weight.
    The remaining factors satisfy
    $\prod_i e^{-t}=e^{-Nt}$ and $\prod_i t^{n_i}=t^m$.
    All factorials are positive, so this argument also applies at $t=0$.
\end{proof}

\begin{theorem}[Independent Poisson count law]
    \label{thm:poisson_independent_counts}
    \lean{PoissonWord.map_countVector, PoissonWord.map_countVector_apply,
        PoissonWord.iIndepFun_countVector}
    \leanok
    \uses{def:poisson_word_measure, def:poisson_word_count_vector}
    Under $\mu_t$, the count vector $w\mapsto(n_i(w))_{i\in I}$ has law
    $\bigotimes_{i\in I}\operatorname{Poisson}(t)$.
    Thus the label counts are mutually independent, and each has
    Poisson law of rate $t$.
\end{theorem}

\begin{proof}
    \leanok
    \uses{lem:poisson_count_event_mass}
    The set $\N^I$ is countable.  Formula~(\ref{eq:poisson_count_event_mass})
    equates the singleton masses of the count law and the product of the
    Poisson laws, so the two measures agree.  Coordinate projection gives
    each marginal law.  Equality of the joint law with the product of its
    marginals proves mutual independence.
\end{proof}

An empty alphabet gives the unique empty count vector and an empty product
of probability measures.  No event times or independent-increment process
are constructed by this fixed-time characterization.
